\section{Ceilings, analytic obstructions and effective collapse}\label{sec:ceilings}
\begin{proposition}[Coupling-exponent ceilings]\label{prop:critical-ceiling}
Suppose finite-volume gaps at rational $g$ are uniformly computable,
defined eventually in volume, and their inner limits exist and are
positive at every sufficiently small positive rational $g$. For finite
endpoints in \eqref{eq:critical-definition}, the lower, upper and
existing ceilings are $\mathsf S_3,\mathsf P_3,\mathsf D_3$.
If the inner limit has a computable error modulus, or a computable
eventual-exactness threshold, these reduce respectively to
$\mathsf S_2,\mathsf P_2,\mathsf D_2$.
If instead the positive inner gap is uniformly right-c.e., the
ceilings are $\mathsf S_3,\mathsf P_2,\mathsf P_2$.
\end{proposition}
\begin{proof}
An existing inner limit has a uniform $0'$-Cauchy name. Promised
positivity lets that oracle find a rational interval bounded away
from zero and compute its logarithm. The normalization is then
uniformly $\mathsf D_2$. Tail infima and suprema over rational
couplings give the same relativized lower/upper cut arguments as
Proposition~\ref{prop:calculus}: use a rational buffer and quantify
over all rational $g$ in a sufficiently small punctured neighbourhood.
An existing outer limit is computable in the next jump by a
tail-Cauchy search. This proves the first ceilings. Effective inner
convergence computes the gap without the first jump, proving the
second. A positive right-c.e. gap has a \emph{left}-c.e. inverse
logarithm, so the $\mathsf S_1$ row proves the third.
\end{proof}

An infimum over an effectively enumerated adequate family of
computable excitation tests is right-c.e. This observation applies
only after the variational representation is established for the
particular gap definition. The certified upper approximations of
Xu et al. \cite{bulkGap} concern their operator-algebraic bulk gap;
they do not identify it with an open-box, film or commensurate-torus
limit. The one-sided ceilings above are bounds, without an attainment
claim for intrinsic bulk dispersion.

\begin{proposition}[Finite analytic obstruction]\label{prop:rellich}
If a finite-dimensional Hermitian matrix depends real-analytically
on $g$ near $g_c$, and its gap above the entire ground space is
defined to the right, that gap has a nonnegative integer exponent.
\end{proposition}
\begin{proof}
Rellich's theorem gives analytic eigenvalue branches
\cite[Chapter~II, Section~6]{kato}. Two branches either agree
identically or have isolated intersections. Finitely many branches
therefore have a fixed ordering of distinct values on some
$(g_c,g_c+\epsilon)$. The gap there is the difference of two analytic
branches, nonzero on that interval. Its Taylor series starts with
$c(g-g_c)^m$, $c>0$ and integer $m\geq0$. Hence its logarithmic
exponent is $m$.
\end{proof}

A noninteger envelope cannot equal the gap of one fixed finite
analytic system throughout a right neighbourhood. In the ruler it
is the increasing cell size that evades this obstruction. This
argument forces infinitely many effective finite pieces in that
representation; it does not force first-order steps in every
possible infinite-volume construction.

\subsection{Locality and first thermodynamic derivatives}
Cutting a large box into side-$L$ boxes removes interactions of norm
at most a computable constant times the number of cut bonds. The
ground-energy variational inequality yields
\begin{equation}\label{eq:density-modulus}
 |e-e_L|\leq C/L.
\end{equation}
The same argument gives a pressure modulus. For Hermitian $A,B$,
Weyl's eigenvalue bounds imply
\[
 |\log\operatorname{Tr}e^{A+B}-\log\operatorname{Tr}e^A|
 \leq\|B\|.
\]
The partition function factors across a cut tensor product.
Tiling and passing to the limit thus give $|p-p_L|\leq C(\beta,h)/L$
uniformly on bounded computable parameter sets. These are two-sided
brackets; an unshifted finite-box value need not itself be a lower
bound. The modulus describes a box side required for accuracy,
not the algorithm's running time.

For a computable convex pressure as a function of an appropriate
linear field variable, at a computable interior point $u$,
\[
 p'_+(u)=\inf_{t\in\Q_{>0}}\frac{p(u+t)-p(u)}t,\qquad
 p'_-(u)=\sup_{t\in\Q_{>0}}\frac{p(u)-p(u-t)}t,
\]
with $t$ restricted to the domain. The right derivative is
right-c.e. and the left derivative left-c.e. At a promised
differentiability point they agree; search for upper and lower chord
bounds within the requested accuracy. This computes the derivative,
uniformly over computable sequences of promised differentiability
points, without a given modulus of differentiability.

\begin{corollary}[Restricted positive-temperature ceiling]
One further normalized limit formed from pressure values or
differentiable first thermodynamic densities, at computable
sampling and reference data, has lower/upper/existing ceilings
$\mathsf S_2,\mathsf P_2,\mathsf D_2$. A positive right derivative
at a cusp, with inverse-log normalization, instead has ceilings
$\mathsf S_3,\mathsf P_2,\mathsf P_2$.
\end{corollary}
\begin{proof}
The locality and convexity arguments compute the first type of
finite-scale data. The second type is right-c.e.; its inverse
logarithm reverses the side. Apply Proposition~\ref{prop:calculus}.
All logarithms and denominators must be defined with the advertised
positivity, and subtracted reference data must have the stated
computability.
\end{proof}

These are ceilings, not thermal attainment theorems. They do not
cover every correlation, higher response or selected Gibbs-state
observable. A noncomputable critical point is additional input and
is not included in a theorem for computable reference data.

\begin{proposition}[Zero-temperature tangent-deficit exponent]
For an affine finite-range family let $g_c$ be computable and put
\[
 s(g)=e'_+(g_c)(g-g_c)-(e(g)-e(g_c)).
\]
Assume $s(g)>0$ at the computable sampling points tending down to
$g_c$ and that
$2-\alpha=\lim_{g\downarrow g_c}\log s(g)/\log(g-g_c)$ exists
and is finite in the stated rational-coupling sense. Then
$2-\alpha\in\mathsf S_2$ and $\alpha\in\mathsf P_2$.
\end{proposition}
\begin{proof}
Equation~\eqref{eq:density-modulus} computes $e$, and the interaction
norm gives a volume-uniform Lipschitz bound in $g$, including $e(g_c)$.
Ground energy is an infimum of affine functions and is concave.
Its right derivative is the supremum of forward chord slopes, so
it is left-c.e. Concavity gives $s\geq0$, and $s$ is uniformly
left-c.e. Under promised strict positivity, $\log s$ is left-c.e.
Division by the negative computable $\log(g-g_c)$ reverses the
side, giving right-c.e. normalized data. Their existing limit is
in $\mathsf S_2$; subtracting from two gives $\mathsf P_2$.
\end{proof}

This defines a particular singular-energy exponent. It matches the
usual $e_{\rm sing}\sim|g-g_c|^{2-\alpha}$ interpretation where that
singular term dominates the analytic remainder, for example
$1<2-\alpha<2$. It is not an unrestricted theorem about every
definition of a specific-heat exponent. Weyl's gap continuity bound
has a volume factor $O(L^d)$ and supplies no analogous uniform
modulus for the gap. A volume-uniform modulus of dependence on $g$
alone, without an effective approximation mechanism in volume,
must not be substituted for an effective inner-limit modulus.
For example computable rational approximants $a_L\to\Omega_U$
give $F_L(g)=a_L$ with common Lipschitz constant zero, yet their
inner limit is an uncomputable constant. This is a data-level
counterexample; the energy-density argument has the additional
effective size bound \eqref{eq:density-modulus}.

\subsection{Other effective mechanisms and their limits}
Finite-range Lieb--Robinson estimates, with a computable initial
state and effective constants, give a computable infinite-volume
local expectation at every finite computable time \cite{liebRobinson}.
A late-time exponent against a computable reference has the
one-limit ceiling. If the reference is itself an untamed limiting
value and is used inside the outer normalization, a further
dependent limit can restore the next ceiling.

For a nonempty effectively closed subset $K$ of an effectively
compact space and computable continuous $f$, $\min_K f$ is left-c.e.:
certify that $K\subset\{f>q\}$ using effective compactness.
Effective overtness also semidecides $K\cap\{f<q\}\ne\varnothing$,
giving the other side and computability. Ordinary compactness alone
does not suffice. These and locality, convexity and variational
extrema are useful sufficient mechanisms, not an exhaustive list
of all ways physical definitions can simplify.

Higher derivatives require separate control. The convex functions
\[
 p_n(u)=u^2+\frac{1-\cos(nu)}{n^2}
\]
converge effectively and uniformly to $u^2$, but $p_n''(0)=3$
whereas the limit has second derivative two. Thus effective pressure
convergence alone does not justify interchanging second derivatives.
Under the ordinary hypotheses giving a second-order expansion, an
existing second derivative of a computable function is itself a
limit of computable central differences and has a $\mathsf D_2$
ceiling. An additional response exponent can then have a
$\mathsf D_3$ ceiling; no higher level follows merely by writing
the derivative as a difference of one-sided values.

If a film uses $\liminf_L$ instead of an existing inner limit,
the inner gap is $\mathsf S_2$; its inverse-log normalization is
$\mathsf P_2$. The lower/upper/existing ceilings are then
$\mathsf S_3,\mathsf P_4,\mathsf S_3$. No physical sharpness for
this alternative convention follows from Theorem~\ref{thm:film}.
Range bounds likewise restrict an interval without proving either
class preservation or attainment.

\subsection{Frustration-free criteria require their hypotheses}
Gosset and Mozgunov give local gap thresholds and, in their gapless
one-dimensional periodic setting, an $O(n^{-2})$ open-chain upper bound
\cite{gossetMozgunov}. Lemm and Mozgunov give finite-size criteria with
boundaries \cite{lemmMozgunov}; their two-dimensional criterion uses a
minimum over rhomboidal patches, and general boundaries require the
corresponding intersected patch shapes. A square-box exponent alone does
not provide those patch or edge estimates. Masaoka, Soejima and Watanabe
obtain $z\geq2$ under their frustration-free correlation hypotheses,
including the extended correlations treated there
\cite{masaokaSoejimaWatanabe}. Lemm and Lucia prove inverse-square critical scaling, up to logarithmic
factors, without a ground-state correlation assumption
\cite[Theorem~1.4 and Remark~1.5]{lemmLucia}. Their result controls an
infimum over specified finite-domain families for one fixed positive
finite-range frustration-free interaction on graphs with the required
Euclidean embedding bounds. A square-box estimate alone does not give
that domain-family infimum, and periodic boundary conditions require the
additional boundary information stated there. These results do not imply
an unconditional exclusion interval for the square-box exponents of the
present setting.

The specified uniform local decomposition here is frustrated: an onsite
flag reward competes with interior bond penalties. For a particular box,
adding $4LI$ lets one regroup the flag energy into positive interior and
boundary terms. That boundary-dependent regrouping does not identify the
original interaction with a uniformly frustration-free family. No general
necessity of frustration for every small exponent is proved here.


